\section{Rational maps}
\label{section:I.7}

\subsection{Rational maps and rational functions}
\label{subsection:I.7.1}

\begin{env}[7.1.1]
\label{I.7.1.1}
Let $X$, $Y$ be two preschemes, $U$, $V$ two dense open subsets of $X$, $f$
(resp. $g$) a morphism from $U$ (resp. $V$) to $Y$; we will say that $f$ and
$g$ are
\emph{equivalent} if they coincide on an open subset dense in $U\cap V$.
Since a finite intersection of dense open subsets of $X$ is a dense open
subset of $X$, it is clear that this relation is an
\emph{equivalence relation}.
\end{env}

\begin{definition}[7.1.2]
\label{I.7.1.2}
Given two preschemes $X$, $Y$, a
\emph{rational map from $X$ to $Y$} is an equivalence class of morphisms from
dense open subsets of $X$ to $Y$, for the relation defined in
(\hyperref[I.7.1.1]{7.1.1}). If $X$ and $Y$ are
$S$-preschemes, we say that such a class is a
\emph{rational $S$-map} if there exists a representative of this class which
is an $S$-morphism. Every rational $S$-map from $S$ to $X$ is called a
\emph{rational $S$-section of $X$}.
Every rational $X$-section of the $X$-prescheme
$X\otimes_{\mathbf{Z}}\mathbf{Z}[T]$ (where $T$ is an indeterminate) is called
a \emph{rational function on a prescheme $X$}.

By abuse of language, when only $S$-preschemes are under consideration, we
will say ``rational map'' instead of ``rational $S$-map'' if no confusion can
result.

Let $f$ be a rational map from $X$ to $Y$, and $U$ an open subset of $X$; if
$f_1$, $f_2$ are morphisms belonging to the class $f$, defined respectively on
dense open subsets $V$, $W$ of $X$, the restrictions $f_1|(U\cap V)$,
$f_2|(U\cap W)$ coincide on an open subset of $U\cap V\cap W$ which is dense
in $U$; the class of morphisms $f$ therefore defines a rational map from $U$
to $Y$, called the
\emph{restriction of $f$ to $U$} and denoted $f|U$.

If to every $S$-morphism $f:X\to Y$ we assign the rational $S$-map
\oldpage[I]{156}
to which $f$ belongs, we define a canonical map from
$\operatorname{Hom}_S(X,Y)$ to the set of rational $S$-maps from $X$ to $Y$.
We denote by $\Gamma_{\mathrm{rat}}(X/Y)$ the set of rational $Y$-sections of
$X$, and we therefore have a canonical map
$\Gamma(X/Y)\to\Gamma_{\mathrm{rat}}(X/Y)$. It is moreover clear that if $X$
and $Y$ are two $S$-preschemes, the set of rational $S$-maps from $X$ to $Y$
is canonically identified with
$\Gamma_{\mathrm{rat}}((X\times_S Y)/X)$
(\hyperref[I.3.3.14]{3.3.14}).
\end{definition}

\begin{env}[7.1.3]
\label{I.7.1.3}
It follows immediately from (\hyperref[I.7.1.2]{7.1.2}) and
(\hyperref[I.3.3.14]{3.3.14}) that the rational functions on $X$ are
canonically identified with the
\emph{equivalence classes of sections of the structural sheaf
$\mathscr{O}_X$} over everywhere dense open subsets of $X$, two such sections
being equivalent if they coincide on an everywhere dense open subset contained
in the intersection of their sets of definition. In particular, it follows
that the rational functions on $X$ form a
\emph{ring} $R(X)$.
\end{env}

\begin{env}[7.1.4]
\label{I.7.1.4}
When $X$ is an \emph{irreducible} prescheme, every non-empty open subset is
dense in $X$; we can also say that the non-empty open subsets of $X$ are the
\emph{open neighbourhoods of the generic point $x$} of $X$. To say that two
morphisms of non-empty open subsets of $X$ into $Y$ are equivalent therefore
means in this case that they have \emph{the same germ} at the point $x$. In
other words, the rational maps (resp. rational $S$-maps)
$X\to Y$ are identified with the \emph{germs of morphisms} (resp. of
$S$-morphisms) of non-empty open subsets of $X$ into $Y$ at the generic point
$x$ of $X$. In particular:
\end{env}

\begin{proposition}[7.1.5]
\label{I.7.1.5}
If $X$ is an irreducible prescheme, the ring $R(X)$ of rational functions
on $X$ is canonically identified with the local ring $\mathscr{O}_x$ of the
generic point $x$ of $X$. It is a local ring of dimension 0, and consequently
an Artinian local ring when $X$ is Noetherian; it is a field when $X$ is
integral, and it is identified with the field of fractions of $A(X)$ when
moreover $X$ is an affine scheme.
\end{proposition}

In view of the above and of the identification of rational functions with
sections of $\mathscr{O}_X$ over an everywhere dense open subset, the first
assertion is nothing but the definition of the stalk of a sheaf at a point.
For the other assertions, we may restrict to the case where $X$ is affine with
ring $A$; then $\mathfrak{j}_x$ is the nilradical of $A$, and $\mathscr{O}_x$
is therefore of dimension 0; if $A$ is integral, $\mathfrak{j}_x=(0)$, and
$\mathscr{O}_x$ is therefore the field of fractions of $A$. Finally, if $A$ is
Noetherian, we know (\cite{I-11}, p.~127, cor.~4) that $\mathfrak{j}_x$ is
nilpotent and $\mathscr{O}_x=A_{\mathfrak{j}_x}$ Artinian.

If $X$ is \emph{integral}, the ring $\mathscr{O}_z$ is integral for every
$z\in X$; every affine open subset $U$ containing $z$ also contains $x$, and
$R(U)$, equal to the field of fractions of $A(U)$, is therefore identified
with $R(X)$; we conclude that $R(X)$ is also identified with the
\emph{field of fractions of $\mathscr{O}_z$}: the canonical identification of
$\mathscr{O}_z$ with a subring of $R(X)$ consists in assigning to every germ of
section $s\in\mathscr{O}_z$ the unique rational function on $X$ that is the
class of a section of $\mathscr{O}_X$ (necessarily defined on an everywhere
dense open subset) having $s$ as its germ at the point $z$.

\begin{env}[7.1.6]
\label{I.7.1.6}
Suppose now that $X$ has a \emph{finite} number of irreducible components
$X_i$ ($1\leq i\leq n$) (which is the case when the underlying space of $X$
is \emph{Noetherian}); let $X'_i$ be the open subset of $X$ which is the
complement, relative to $X_i$, of the union of the $X_j\cap X_i$ for
$j\neq i$; $X'_i$ is irreducible, its generic point $x_i$ is the generic point
of $X_i$, and the $X'_i$ are pairwise disjoint, their union being dense in $X$
(\hyperref[0.2.1.6]{0, 2.1.6}). For
\oldpage[I]{157}
each everywhere dense open subset $U$ of $X$, $U_i=U\cap X'_i$ is a non-empty
open subset dense in $X'_i$, the $U_i$ having pairwise no point in common,
hence $U'=\bigcup_i U_i$ is dense in $X$. To give a morphism from $U'$ to $Y$
amounts to giving (arbitrarily) a morphism from each of the $U_i$ to $Y$.
Hence:
\end{env}

\begin{proposition}[7.1.7]
\label{I.7.1.7}
Let $X$, $Y$ be two preschemes (resp. $S$-preschemes) such that $X$ has a
finite number of irreducible components $X_i$, with generic points $x_i$
($1\leq i\leq n$). If $R_i$ is the set of germs of morphisms (resp.
$S$-morphisms) of open subsets of $X$ into $Y$ at the point $x_i$, the set of
rational maps (resp. rational $S$-maps) from $X$ to $Y$ is identified with the
product of the $R_i$ ($1\leq i\leq n$).
\end{proposition}

\begin{corollary}[7.1.8]
\label{I.7.1.8}
Let $X$ be a Noetherian prescheme. The ring of rational functions on $X$ is
an Artinian ring, whose local components are the rings
$\mathscr{O}_{x_i}$ of the generic points $x_i$ of the irreducible components
of $X$.
\end{corollary}

\begin{corollary}[7.1.9]
\label{I.7.1.9}
Let $A$ be a Noetherian ring, and let $X=\operatorname{Spec}(A)$. If $Q$ is
the complement of the union of the minimal prime ideals of $A$, the ring of
rational functions on $X$ is canonically identified with the ring of
fractions $Q^{-1}A$.
\end{corollary}

This will follow from the following lemma:

\begin{lemma}[7.1.9.1]
\label{I.7.1.9.1}
For an element $f\in A$ to be such that $D(f)$ is everywhere dense in $X$, it
is necessary and sufficient that $f\in Q$; every dense open subset of $X$
contains an open subset of the form $D(f)$, where $f\in Q$.
\end{lemma}

Indeed, suppose this lemma proved; since the ring of sections
$\Gamma(D(f),\mathscr{O}_X)$ is identified with $A_f$
(\hyperref[I.1.3.6]{1.3.6} and \hyperref[I.1.3.7]{1.3.7}), it follows
from the fact that the $D(f)$ with $f\in Q$ form a cofinal set in the set
(ordered by $\supset$) of dense open subsets of $X$, and from the definition
(\hyperref[I.7.1.1]{7.1.1}), that the ring of rational functions on $X$ is
identified with the inductive limit of the $A_f$ for $f\in Q$ (for the
preorder relation ``$g$ is a multiple of $f$''), that is, with $Q^{-1}A$
(\hyperref[0.1.4.5]{0, 1.4.5}).

To prove (\hyperref[I.7.1.9.1]{7.1.9.1}), again denote by $X_i$ the
irreducible components of $X$ ($1\leq i\leq n$); if $D(f)$ is dense in $X$,
then $D(f)\cap X_i\neq\varnothing$ for $1\leq i\leq n$, and conversely; but
this means that $f\notin\mathfrak{p}_i$ for $1\leq i\leq n$, setting
$\mathfrak{p}_i=\mathfrak{j}(X_i)$, and since the $\mathfrak{p}_i$ are the
minimal prime ideals of $A$ (\hyperref[I.1.1.14]{1.1.14}), the relations
$f\notin\mathfrak{p}_i$ ($1\leq i\leq n$) are equivalent to $f\in Q$, whence
the first assertion of the lemma. On the other hand, if $U$ is a dense open
subset of $X$, the complement of $U$ is a set of the form $V(\mathfrak{a})$,
where $\mathfrak{a}$ is an ideal which is contained in none of the
$\mathfrak{p}_i$; it is therefore not contained in their union (\cite{I-10},
p.~13), and so there exists an $f\in\mathfrak{a}$ belonging to $Q$; whence
$D(f)\subset U$, which completes the proof.

\begin{env}[7.1.10]
\label{I.7.1.10}
Suppose again that $X$ is irreducible, with generic point $x$. Since every
non-empty open subset $U$ of $X$ contains $x$, and consequently also contains
every $z\in X$ such that $x\in\overline{\{z\}}$, every morphism $U\to Y$ can
be composed with the canonical morphism $\operatorname{Spec}(\mathscr{O}_x)\to X$
(\hyperref[I.2.4.1]{2.4.1}); and two morphisms into $Y$ of two non-empty
open subsets of $X$ which coincide on a non-empty open subset of $X$ give, by
composition, the same morphism $\operatorname{Spec}(\mathscr{O}_x)\to Y$. In
other words, to every rational map from $X$ to $Y$ there thus corresponds a
well-determined morphism $\operatorname{Spec}(\mathscr{O}_x)\to Y$.
\end{env}

\oldpage[I]{158}
\begin{proposition}[7.1.11]
\label{I.7.1.11}
Let $X$, $Y$ be two $S$-preschemes; suppose $X$ irreducible, with generic
point $x$, and $Y$ of finite type over $S$. Two rational $S$-maps from $X$ to
$Y$ to which the same $S$-morphism $\operatorname{Spec}(\mathscr{O}_x)\to Y$
corresponds are identical. If we suppose moreover that $S$ is locally
Noetherian, every $S$-morphism from $\operatorname{Spec}(\mathscr{O}_x)$ to
$Y$ corresponds to a rational $S$-map (and only one) from $X$ to $Y$.
\end{proposition}

Taking into account the fact that every non-empty open subset of $X$ is
everywhere dense, this follows immediately from
(\hyperref[I.6.5.1]{6.5.1}).

\begin{corollary}[7.1.12]
\label{I.7.1.12}
Suppose $S$ locally Noetherian, and the other hypotheses of
(\hyperref[I.7.1.11]{7.1.11}) satisfied. The rational $S$-maps from $X$ to
$Y$ are then identified with the points of the $S$-prescheme $Y$ with values
in the $S$-prescheme $\operatorname{Spec}(\mathscr{O}_x)$.
\end{corollary}

This is nothing but (\hyperref[I.7.1.11]{7.1.11}), with the terminology
introduced in (\hyperref[I.3.4.1]{3.4.1}).

\begin{corollary}[7.1.13]
\label{I.7.1.13}
Suppose the conditions of (\hyperref[I.7.1.12]{7.1.12}) are fulfilled. Let
$s$ be the image of $x$ in $S$. To give a rational $S$-map from $X$ to $Y$ is
equivalent to giving a point $y$ of $Y$ above $s$ and a local
$\mathscr{O}_s$-homomorphism
$\mathscr{O}_y\to\mathscr{O}_x=R(X)$.
\end{corollary}

This follows from (\hyperref[I.7.1.11]{7.1.11}) and
(\hyperref[I.2.4.4]{2.4.4}).

In particular:

\begin{corollary}[7.1.14]
\label{I.7.1.14}
Under the conditions of (\hyperref[I.7.1.12]{7.1.12}), the rational
$S$-maps from $X$ to $Y$ depend (for given $Y$) only on the
$S$-prescheme $\operatorname{Spec}(\mathscr{O}_x)$, and in particular remain
the same when $X$ is replaced by
$\operatorname{Spec}(\mathscr{O}_z)$, for any $z\in X$.
\end{corollary}

Indeed, since $z\in\overline{\{x\}}$, $x$ is the generic point of
$Z=\operatorname{Spec}(\mathscr{O}_z)$, and
$\mathscr{O}_{X,x}=\mathscr{O}_{Z,x}$.

When $X$ is integral, $R(X)=\mathscr{O}_x=k(x)$ is a field
(\hyperref[I.7.1.5]{7.1.5}); the preceding corollaries then specialize
to:

\begin{corollary}[7.1.15]
\label{I.7.1.15}
Suppose the conditions of (\hyperref[I.7.1.12]{7.1.12}) are satisfied and
moreover that $X$ is integral. Let $s$ be the image of $x$ in $S$. Then the
rational $S$-maps from $X$ to $Y$ are identified with the geometric points of
$Y\otimes_S k(s)$ with values in the extension $R(X)$ of $k(s)$; in other
words, each of them is equivalent to giving a point $y\in Y$ above $s$ and a
$k(s)$-monomorphism from $k(y)$ to
$k(x)=R(X)$.
\end{corollary}

Indeed, the points of $Y$ above $s$ are identified with those of
$Y\otimes_S k(s)$ (\hyperref[I.3.6.3]{3.6.3}), and the local
$\mathscr{O}_s$-homomorphisms $\mathscr{O}_y\to R(X)$ with the
$k(s)$-monomorphisms $k(y)\to R(X)$.

More particularly:

\begin{corollary}[7.1.16]
\label{I.7.1.16}
Let $k$ be a field, $X$, $Y$ two algebraic preschemes
(\hyperref[I.6.4.1]{6.4.1}) over $k$; suppose moreover that $X$ is
integral. Then the rational $k$-maps from $X$ to $Y$ are identified with the
geometric points of $Y$ with values in the extension $R(X)$ of $k$
(\hyperref[I.3.4.4]{3.4.4}).
\end{corollary}

\subsection{Domain of definition of a rational map}
\label{subsection:I.7.2}

\begin{env}[7.2.1]
\label{I.7.2.1}
Let $X$, $Y$ be two preschemes, $f$ a rational map from $X$ to $Y$. We say
that $f$ is \emph{defined at a point $x\in X$} if there exist an everywhere
dense open set $U$ containing $x$ and a morphism $U\to Y$ belonging to the
equivalence class $f$. The set of points $x\in X$ where $f$ is defined is
called the \emph{domain of definition of $f$}; it is clear that it is an
everywhere dense open subset of $X$.
\end{env}

\oldpage[I]{159}
\begin{proposition}[7.2.2]
\label{I.7.2.2}
Let $X$, $Y$ be two $S$-preschemes such that $X$ is reduced and $Y$ is
separated over $S$. Let $f$ be a rational $S$-map from $X$ to $Y$, and $U_0$
its domain of definition. Then there exists one and only one $S$-morphism
$U_0\to Y$ belonging to the class $f$.
\end{proposition}

Since for every morphism $U\to Y$ belonging to the class $f$ we necessarily
have $U\subset U_0$, it is clear that the proposition will be a consequence
of

\begin{lemma}[7.2.2.1]
\label{I.7.2.2.1}
Under the hypotheses of (\hyperref[I.7.2.2]{7.2.2}), let $U_1$, $U_2$ be
two everywhere dense open subsets of $X$, and $f_i:U_i\to Y$ ($i=1,2$) two
$S$-morphisms such that there exists an open subset $V\subset U_1\cap U_2$, dense
in $X$, on which $f_1$ and $f_2$ coincide. Then $f_1$ and $f_2$ coincide on
$U_1\cap U_2$.
\end{lemma}

We may evidently restrict to the case where $X=U_1=U_2$. Since $X$ (hence
$V$) is reduced, $X$ is the smallest closed sub-prescheme of $X$ dominating $V$
(\hyperref[I.5.2.2]{5.2.2}). Let
$g=(f_1,f_2)_S:X\to Y\times_S Y$; since by hypothesis the diagonal
$T=\Delta_Y(Y)$ is a closed sub-prescheme of $Y\times_S Y$,
$Z=g^{-1}(T)$ is a closed sub-prescheme of $X$
(\hyperref[I.4.4.1]{4.4.1}). If $h:V\to Y$ is the common restriction of
$f_1$ and $f_2$ to $V$, the restriction of $g$ to $V$ is $g'=(h,h)_S$, which
factors as $g'=\Delta_Y\circ h$; since $\Delta_Y^{-1}(T)=Y$, we have
$g'^{-1}(T)=V$, and consequently $Z$ is a closed sub-prescheme of $X$ inducing
$V$, hence dominating $V$, which implies $Z=X$. From the relation
$g^{-1}(T)=X$ we deduce (\hyperref[I.4.4.1]{4.4.1}) that $g$ factors as
$\Delta_Y\circ f$, where $f$ is a morphism $X\to Y$, which implies, by the
definition of the diagonal morphism, that $f_1=f_2=f$.

It is clear that the morphism $U_0\to Y$ defined in
(\hyperref[I.7.2.2]{7.2.2}) is the unique morphism of the class $f$ which
\emph{cannot be extended} to a morphism of an open subset of $X$ strictly
containing $U_0$. \emph{Under the hypotheses of
(\hyperref[I.7.2.2]{7.2.2})}, we can therefore \emph{identify} the
rational maps from $X$ to $Y$ with the \emph{non-extendable} morphisms (to
strictly larger open subsets) of everywhere dense open subsets of $X$ into $Y$.
With this identification, Proposition
(\hyperref[I.7.2.2]{7.2.2}) implies:

\begin{corollary}[7.2.3]
\label{I.7.2.3}
The hypotheses on $X$ and $Y$ being those of
(\hyperref[I.7.2.2]{7.2.2}), let $U$ be an everywhere dense open subset of
$X$. There exists a canonical bijective correspondence between the
$S$-morphisms from $U$ to $Y$ and the rational $S$-maps from $X$ to $Y$
defined at all points of $U$.
\end{corollary}

Indeed, by virtue of (\hyperref[I.7.2.2]{7.2.2}), for every $S$-morphism
$f$ from $U$ to $Y$ there exists one and only one rational $S$-map
$\overline f$ from $X$ to $Y$ which extends $f$.

\begin{corollary}[7.2.4]
\label{I.7.2.4}
Let $S$ be a scheme, $X$ a reduced $S$-prescheme, $Y$ an $S$-scheme, and
$f:U\to Y$ an $S$-morphism of a dense open subset $U$ of $X$ into $Y$. If
$\overline f$ is the rational $\mathbf{Z}$-map from $X$ to $Y$ which
extends $f$, then $\overline f$ is an $S$-morphism (and is consequently the
rational $S$-map from $X$ to $Y$ extending $f$).
\end{corollary}

Indeed, if $\varphi:X\to S$, $\psi:Y\to S$ are the structural morphisms,
$U_0$ the domain of definition of $\overline f$, and $j$ the injection
$U_0\to X$, it suffices to prove that $\psi\circ\overline f=\varphi\circ j$,
which follows immediately from (\hyperref[I.7.2.2.1]{7.2.2.1}), since $f$
is an $S$-morphism.

\begin{corollary}[7.2.5]
\label{I.7.2.5}
Let $X$, $Y$ be two $S$-preschemes; suppose $X$ reduced, and $X$ and $Y$
separated over $S$. Let $p:Y\to X$ be an $S$-morphism (making $Y$ an
$X$-prescheme), $U$ an everywhere dense open subset of $X$, and $f$ a
$U$-section of $Y$; then the rational map $\overline f$ from $X$ to $Y$
extending $f$ is a rational $X$-section of $Y$.
\end{corollary}

\oldpage[I]{160}
We must prove that $p\circ\overline f$ is the identity on the domain of
definition of $\overline f$; since $X$ is separated over $S$, this again
follows from (\hyperref[I.7.2.2.1]{7.2.2.1}).

\begin{corollary}[7.2.6]
\label{I.7.2.6}
Let $X$ be a reduced prescheme, and $U$ an everywhere dense open subset of
$X$. There is a canonical bijective correspondence between the sections of
$\mathscr{O}_X$ over $U$ and the rational functions $f$ on $X$ defined at
every point of $U$.
\end{corollary}

Taking into account (\hyperref[I.7.2.3]{7.2.3}),
(\hyperref[I.7.1.2]{7.1.2}) and (\hyperref[I.7.1.3]{7.1.3}), it
suffices to note that the $X$-prescheme $X\otimes_{\mathbf Z}\mathbf Z[T]$ is
separated over $X$ (\hyperref[I.5.5.1]{5.5.1, (iv)}).

\begin{corollary}[7.2.7]
\label{I.7.2.7}
Let $Y$ be a reduced prescheme, $f:X\to Y$ a separated morphism, $U$ an
everywhere dense open subset of $Y$, $g:U\to f^{-1}(U)$ a $U$-section of
$f^{-1}(U)$, and $Z$ the reduced sub-prescheme of $X$ having
$\overline{g(U)}$ as underlying space
(\hyperref[I.5.2.1]{5.2.1}). For $g$ to be the restriction of a
$Y$-section of $X$ (in other words (\hyperref[I.7.2.5]{7.2.5}), for the
rational map from $Y$ to $X$ extending $g$ to be everywhere defined), it is
necessary and sufficient that the restriction of $f$ to $Z$ be an
isomorphism of $Z$ onto $Y$.
\end{corollary}

The restriction of $f$ to $f^{-1}(U)$ is a separated morphism
(\hyperref[I.5.5.1]{5.5.1, (i)}), hence $g$ is a closed immersion
(\hyperref[I.5.4.6]{5.4.6}), and consequently
$g(U)=Z\cap f^{-1}(U)$ and the sub-prescheme induced by $Z$ on the open subset
$g(U)$ of $Z$ is identical to the closed sub-prescheme of $f^{-1}(U)$
associated to $g$ (\hyperref[I.5.2.1]{5.2.1}). It is then clear that the
condition of the statement is sufficient, since if it is fulfilled, and if
$f_Z:Z\to Y$ is the restriction of $f$ to $Z$ and $\overline g:Y\to Z$ the
inverse isomorphism, then $\overline g$ extends $g$. Conversely, if $g$ is the
restriction to $U$ of a $Y$-section $h$ of $X$, then $h$ is a closed immersion
(\hyperref[I.5.4.6]{5.4.6}), hence $h(Y)$ is closed, and since it is
contained in $Z$, it is equal to $Z$, and it follows from
(\hyperref[I.5.2.1]{5.2.1}) that $h$ is necessarily an isomorphism of
$Y$ onto the closed sub-prescheme $Z$ of $X$.

\begin{env}[7.2.8]
\label{I.7.2.8}
Let $X$, $Y$ be two $S$-preschemes, $X$ being assumed reduced and $Y$
separated over $S$. Let $f$ be a rational $S$-map from $X$ to $Y$, and let
$x$ be a point of $X$; we can compose $f$ with the canonical $S$-morphism
$\operatorname{Spec}(\mathscr{O}_x)\to X$
(\hyperref[I.2.4.1]{2.4.1}) provided that the trace on
$\operatorname{Spec}(\mathscr{O}_x)$ of the domain of definition of $f$ is
dense in $\operatorname{Spec}(\mathscr{O}_x)$ (identified with the set of
$z\in X$ such that $x\in\overline{\{z\}}$
(\hyperref[I.2.4.2]{2.4.2})). This will be so in the following cases:
\begin{enumerate}
  \item[1\textsuperscript{o}] $X$ is irreducible (hence integral), since then
    the generic point $\xi$ of $X$ is the generic point of
    $\operatorname{Spec}(\mathscr{O}_x)$; since the domain of definition $U$
    of $f$ contains $\xi$, $U\cap\operatorname{Spec}(\mathscr{O}_x)$ contains
    $\xi$, hence is dense in $\operatorname{Spec}(\mathscr{O}_x)$.
  \item[2\textsuperscript{o}] $X$ is locally Noetherian; indeed, our
    assertion then follows from
\end{enumerate}
\end{env}

\begin{lemma}[7.2.8.1]
\label{I.7.2.8.1}
Let $X$ be a prescheme whose underlying space is locally Noetherian, and
$x$ a point of $X$. The irreducible components of
$\operatorname{Spec}(\mathscr{O}_x)$ are the traces on
$\operatorname{Spec}(\mathscr{O}_x)$ of the irreducible components of $X$
containing $x$. For an open subset $U\subset X$ to be such that
$U\cap\operatorname{Spec}(\mathscr{O}_x)$ is dense in
$\operatorname{Spec}(\mathscr{O}_x)$, it is necessary and sufficient that it
meet the irreducible components of $X$ containing $x$ (which is the case in
particular if $U$ is dense in $X$).
\end{lemma}

The second assertion evidently follows from the first, and it therefore
suffices to prove the latter. Since $\operatorname{Spec}(\mathscr{O}_x)$ is
contained in every affine open subset $U$ containing $x$, and the irreducible
components of $U$ containing $x$ are the traces on $U$ of the irreducible
components of $X$ containing $x$
(\hyperref[0.2.1.6]{0, 2.1.6}), we may suppose $X$ affine with ring $A$.
Since the prime ideals of $A_x$ correspond bijectively to the prime ideals
\oldpage[I]{161}
of $A$ contained in $\mathfrak{j}_x$
(\hyperref[0.1.2.6]{0, 1.2.6}), the minimal prime ideals of $A_x$
correspond to the minimal prime ideals of $A$ contained in
$\mathfrak{j}_x$, whence the lemma (\hyperref[I.1.1.14]{1.1.14}).

This being so, suppose we are in one of the two cases mentioned above. If $U$
is the domain of definition of the rational $S$-map $f$, denote by $f'$ the
rational map from $\operatorname{Spec}(\mathscr{O}_x)$ to $Y$ which coincides
(taking into account (\hyperref[I.2.4.2]{2.4.2})) with $f$ on
$U\cap\operatorname{Spec}(\mathscr{O}_x)$; we will say that this rational map
is \emph{induced by $f$}.

\begin{proposition}[7.2.9]
\label{I.7.2.9}
Let $S$ be a locally Noetherian prescheme, $X$ a reduced $S$-prescheme, and
$Y$ an $S$-scheme of finite type. Suppose moreover that $X$ is irreducible or
locally Noetherian. Let then $f$ be a rational $S$-map from $X$ to $Y$, and
$x$ a point of $X$. For $f$ to be defined at the point $x$, it is necessary
and sufficient that the rational map $f'$ from
$\operatorname{Spec}(\mathscr{O}_x)$ to $Y$ induced by $f$
(\hyperref[I.7.2.8]{7.2.8}) be a morphism.
\end{proposition}

The condition being evidently necessary (since
$\operatorname{Spec}(\mathscr{O}_x)$ is contained in every open subset containing
$x$), let us prove that it is sufficient. By virtue of
(\hyperref[I.6.5.1]{6.5.1}), there exist an open neighbourhood $U$ of $x$
in $X$ and an $S$-morphism $g$ from $U$ to $Y$ inducing $f'$ on
$\operatorname{Spec}(\mathscr{O}_x)$. If $X$ is irreducible, $U$ is dense in
$X$, and by virtue of (\hyperref[I.7.2.3]{7.2.3}) we may suppose that $g$
is a rational $S$-map. Moreover, the generic point of $X$ belongs to
$\operatorname{Spec}(\mathscr{O}_x)$ and to the domain of definition of $f$,
hence $f$ and $g$ coincide at this point, and consequently on a non-empty open
set of $X$ (\hyperref[I.6.5.1]{6.5.1}). But since $f$ and $g$ are rational
$S$-maps, they are identical (\hyperref[I.7.2.3]{7.2.3}), hence $f$ is
defined at $x$.

If we now suppose $X$ locally Noetherian, we may suppose $U$ Noetherian;
there are then only finitely many irreducible components $X_i$ of $X$
containing $x$ (\hyperref[I.7.2.8.1]{7.2.8.1}), and we may suppose that
they are the only ones meeting $U$, replacing $U$ if necessary by a smaller
open subset (since there are only finitely many irreducible components of $X$
meeting $U$, $U$ being Noetherian). We then see as above that $f$ and $g$
coincide on a non-empty open subset of each of the $X_i$. Taking into account
the fact that each of the $X_i$ is contained in $\overline U$, consider then
the morphism $f_1$, defined on a dense open subset of $U\cup(X-\overline U)$,
equal to $g$ on $U$ and to $f$ on the intersection of $X-\overline U$ and the
domain of definition of $f$. Since $U\cup(X-\overline U)$ is dense in $X$,
$f_1$ and $f$ coincide on a dense open subset of $X$, and since $f$ is a
rational map, $f$ is an extension of $f_1$ (\hyperref[I.7.2.3]{7.2.3}),
hence is defined at the point $x$.

\subsection{Sheaf of rational functions}
\label{subsection:I.7.3}

\begin{env}[7.3.1]
\label{I.7.3.1}
Let $X$ be a prescheme. For every open subset $U\subset X$, denote by $R(U)$
the ring of rational functions on $U$
(\hyperref[I.7.1.3]{7.1.3}); it is a
$\Gamma(U,\mathscr{O}_X)$-algebra. Moreover, if $V\subset U$ is a second
open subset of $X$, every section of $\mathscr{O}_X$ over an everywhere dense
open subset of $U$ gives, by restriction to $V$, a section over an everywhere
dense open subset of $V$, and if two sections coincide over an everywhere
dense open subset of $U$, their restrictions to $V$ coincide over an
everywhere dense open subset of $V$. We thus define a di-homomorphism of
algebras $\rho_{V,U}:R(U)\to R(V)$, and it is clear that if
$U\supset V\supset W$ are three open subsets of $X$, we have
$\rho_{W,U}=\rho_{W,V}\circ\rho_{V,U}$; the $R(U)$ therefore define a
\emph{presheaf} of algebras on $X$.
\end{env}

\oldpage[I]{162}
\begin{definition}[7.3.2]
\label{I.7.3.2}
The \emph{sheaf of rational functions} on a prescheme $X$, denoted by
$\mathscr{R}(X)$, is the $\mathscr{O}_X$-algebra associated to the presheaf
formed by the $R(U)$.
\end{definition}

For every prescheme $X$ and every open subset $U\subset X$, it is clear that the
induced sheaf $\mathscr{R}(X)|U$ is nothing but $\mathscr{R}(U)$.

\begin{proposition}[7.3.3]
\label{I.7.3.3}
Let $X$ be a prescheme such that the family $(X_\lambda)$ of its irreducible
components is locally finite (which is in particular the case when the
underlying space of $X$ is locally Noetherian). Then the
$\mathscr{O}_X$-module $\mathscr{R}(X)$ is quasi-coherent, and for every open
subset $U$ of $X$ meeting only finitely many components $X_\lambda$,
$R(U)$ is equal to $\Gamma(U,\mathscr{R}(X))$ and is canonically identified
with the direct product of the local rings of the generic points of the
$X_\lambda$ such that $U\cap X_\lambda\neq\varnothing$.
\end{proposition}

We may evidently restrict to the case where $X$ has only finitely many
irreducible components $X_i$, with generic points $x_i$
($1\leq i\leq n$). The fact that $R(U)$ is canonically identified with the
direct product of the $\mathscr{O}_{x_i}=R(X_i)$ such that
$U\cap X_i\neq\varnothing$ then follows from
(\hyperref[I.7.1.7]{7.1.7}). Let us show moreover that the presheaf
$U\to R(U)$ satisfies the sheaf axioms, which will prove that
$R(U)=\Gamma(U,\mathscr{R}(X))$. Indeed, it satisfies (F 1) by the above. To
see that it satisfies (F 2), consider a covering of an open subset $U$ of $X$
by open subsets $V_\alpha\subset U$; if the $s_\alpha\in R(V_\alpha)$ are such
that the restrictions of $s_\alpha$ and $s_\beta$ to $V_\alpha\cap V_\beta$
coincide for every pair of indices, we conclude that for every index $i$ such
that $U\cap X_i\neq\varnothing$, the components in $R(X_i)$ of all the
$s_\alpha$ such that $V_\alpha\cap X_i\neq\varnothing$ are the same; denoting
this component by $t_i$, it is clear that the element of $R(U)$ having the
$t_i$ as components has restriction $s_\alpha$ to each $V_\alpha$. Finally, to
see that $\mathscr{R}(X)$ is quasi-coherent, we may restrict to the case where
$X=\operatorname{Spec}(A)$ is affine; taking for $U$ the affine open subsets
of the form $D(f)$, where $f\in A$, it follows from the above and from the
definition (\hyperref[I.1.3.4]{1.3.4}) that
$\mathscr{R}(X)=\widetilde M$, where $M$ is the direct sum of the $A$-modules
$A_{x_i}$.

\begin{corollary}[7.3.4]
\label{I.7.3.4}
Let $X$ be a reduced prescheme having only finitely many irreducible
components, and let $X_i$ ($1\leq i\leq n$) be the reduced closed
sub-preschemes of $X$ having the irreducible components of $X$ as underlying
spaces (\hyperref[I.5.2.1]{5.2.1}). If $h_i$ is the canonical injection
$X_i\to X$, then $\mathscr{R}(X)$ is the direct product of the
$\mathscr{O}_X$-algebras $(h_i)_*(\mathscr{R}(X_i))$.
\end{corollary}

\begin{corollary}[7.3.5]
\label{I.7.3.5}
If $X$ is irreducible, every quasi-coherent $\mathscr{R}(X)$-module
$\mathscr{F}$ is a simple sheaf.
\end{corollary}

It suffices to show that every $x\in X$ admits a neighbourhood $U$ such that
$\mathscr{F}|U$ is a simple sheaf
(\hyperref[0.3.6.2]{0, 3.6.2}); in other words, we are reduced to the case
where $X$ is affine; we may moreover suppose that $\mathscr{F}$ is the
cokernel of a homomorphism
$(\mathscr{R}(X))^{(I)}\to(\mathscr{R}(X))^{(J)}$
(\hyperref[0.5.1.3]{0, 5.1.3}), and everything comes down to seeing that
$\mathscr{R}(X)$ is a simple sheaf; but this is evident since
$\Gamma(U,\mathscr{R}(X))=R(X)$ for every non-empty open subset $U$, $U$
containing the generic point of $X$.

\begin{corollary}[7.3.6]
\label{I.7.3.6}
If $X$ is irreducible, then for every quasi-coherent $\mathscr{O}_X$-module
$\mathscr{F}$, $\mathscr{F}\otimes_{\mathscr{O}_X}\mathscr{R}(X)$ is a simple
sheaf; if moreover $X$ is reduced (hence integral),
$\mathscr{F}\otimes_{\mathscr{O}_X}\mathscr{R}(X)$ is isomorphic to a sheaf of
the form $(\mathscr{R}(X))^{(I)}$.
\end{corollary}

The second assertion follows from the fact that $R(X)$ is then a field.

\begin{proposition}[7.3.7]
\label{I.7.3.7}
Suppose that the prescheme $X$ is locally integral or locally
\oldpage[I]{163}
Noetherian. Then $\mathscr{R}(X)$ is a quasi-coherent $\mathscr{O}_X$-algebra;
if moreover $X$ is reduced (which is the case when $X$ is locally integral),
the canonical homomorphism $\mathscr{O}_X\to\mathscr{R}(X)$ is injective.
\end{proposition}

The question being local, the first assertion follows from
(\hyperref[I.7.3.3]{7.3.3}); the second follows immediately from
(\hyperref[I.7.2.3]{7.2.3}).

\begin{env}[7.3.8]
\label{I.7.3.8}
Let $X,Y$ be two integral preschemes, which implies that $\mathscr{R}(X)$
(resp. $\mathscr{R}(Y)$) is a quasi-coherent $\mathscr{O}_X$-module (resp.
$\mathscr{O}_Y$-module) (\hyperref[I.7.3.3]{7.3.3}). Let $f:X\to Y$ be a
dominant morphism; then there exists a canonical homomorphism of
$\mathscr{O}_X$-modules
\[
  \tau:f^*(\mathscr{R}(Y))\to\mathscr{R}(X).
\tag{7.3.8.1}
\label{I.7.3.8.1}
\]
\end{env}

First suppose $X=\operatorname{Spec}(A)$ and $Y=\operatorname{Spec}(B)$
affine with integral rings $A$ and $B$, $f$ thus corresponding to an injective
homomorphism $B\to A$, which extends to a monomorphism $L\to K$ from the field
of fractions $L$ of $B$ into the field of fractions $K$ of $A$. The
homomorphism (7.3.8.1) then corresponds to the canonical homomorphism
\[
  L\otimes_B A\to K
\]
(\hyperref[I.1.6.5]{1.6.5}). In the general case, for every pair of
non-empty affine open subsets $U\subset X$, $V\subset Y$ such that
$f(U)\subset V$, we define a homomorphism $\tau_{U,V}$ in the preceding
manner, and we see immediately that if $U'\subset U$, $V'\subset V$,
$f(U')\subset V'$, then $\tau_{U,V}$ extends $\tau_{U',V'}$, whence our
assertion. If $x,y$ are the generic points of $X$ and $Y$ respectively, we
have $f(x)=y$,
\[
  \bigl(f^*(\mathscr{R}(Y))\bigr)_x
  =\mathscr{O}_y\otimes_{\mathscr{O}_y}\mathscr{O}_x
  =\mathscr{O}_x
\]
(\hyperref[0.4.3.1]{0, 4.3.1}), and $\tau_x$ is therefore an isomorphism.

\subsection{Torsion sheaves and torsion-free sheaves}
\label{subsection:I.7.4}

\begin{env}[7.4.1]
\label{I.7.4.1}
Let $X$ be an \emph{integral} prescheme. For every $\mathscr{O}_X$-module
$\mathscr{F}$, the canonical homomorphism
$\mathscr{O}_X\to\mathscr{R}(X)$ defines, by tensoring, a homomorphism
(also called \emph{canonical})
\[
  \mathscr{F}\to
  \mathscr{F}\otimes_{\mathscr{O}_X}\mathscr{R}(X)
\]
which, on each stalk, is nothing but the homomorphism $z\to z\otimes 1$ from
$\mathscr{F}_x$ to
$\mathscr{F}_x\otimes_{\mathscr{O}_x}R(X)$. The \emph{kernel $\mathscr{T}$}
of this homomorphism is a sub-$\mathscr{O}_X$-module of $\mathscr{F}$,
called the \emph{torsion sheaf of $\mathscr{F}$}; it is quasi-coherent if
$\mathscr{F}$ is quasi-coherent
(\hyperref[I.4.1.1]{4.1.1} and \hyperref[I.7.3.6]{7.3.6}). We say that
$\mathscr{F}$ is \emph{torsion-free} if $\mathscr{T}=0$, and that
$\mathscr{F}$ is a \emph{torsion sheaf} if
$\mathscr{T}=\mathscr{F}$. For every $\mathscr{O}_X$-module $\mathscr{F}$,
$\mathscr{F}/\mathscr{T}$ is torsion-free. We deduce from
(\hyperref[I.7.3.5]{7.3.5}) that:
\end{env}

\begin{proposition}[7.4.2]
\label{I.7.4.2}
If $X$ is an integral prescheme, every torsion-free quasi-coherent
$\mathscr{O}_X$-module $\mathscr{F}$ is isomorphic to a sub-sheaf $\mathscr{G}$
of a simple sheaf of the form $(\mathscr{R}(X))^{(I)}$, generated (as an
$\mathscr{R}(X)$-module) by $\mathscr{G}$.
\end{proposition}

The cardinal of $I$ is called the \emph{rank of $\mathscr{F}$}; for every
non-empty affine open subset $U$ of $X$, the rank of $\mathscr{F}$ is equal
to the rank of $\Gamma(U,\mathscr{F})$ as a $\Gamma(U,\mathscr{O}_X)$-module,
as is seen immediately by considering the generic point of $X$, which is
contained in $U$. In particular:

\begin{corollary}[7.4.3]
\label{I.7.4.3}
On an integral prescheme $X$, every torsion-free quasi-coherent
$\mathscr{O}_X$-module of rank $1$ (in particular every invertible
$\mathscr{O}_X$-module) is isomorphic to a nonzero quasi-coherent
sub-$\mathscr{O}_X$-module of $\mathscr{R}(X)$, and conversely.
\end{corollary}

\begin{corollary}[7.4.4]
\label{I.7.4.4}
Let $X$ be an integral prescheme, $\mathscr{L}$, $\mathscr{L}'$ two
torsion-free $\mathscr{O}_X$-modules, and $f$ (resp. $f'$) a section of
$\mathscr{L}$ (resp. $\mathscr{L}'$) over $X$. For
$f\otimes f'=0$, it is necessary and sufficient that one of the sections $f$,
$f'$ be zero.
\end{corollary}

Let $x$ be the generic point of $X$; by hypothesis we have
$(f\otimes f')_x=f_x\otimes f'_x=0$. Since $\mathscr{L}_x$ and
$\mathscr{L}'_x$ are vector spaces over the field
$\mathscr{O}_x$, the preceding relation implies $f_x=0$ or $f'_x=0$, and
consequently $f=0$ or $f'=0$, since $\mathscr{L}$ and $\mathscr{L}'$ are
torsion-free (\hyperref[I.7.3.5]{7.3.5}).

\begin{proposition}[7.4.5]
\label{I.7.4.5}
Let $X$, $Y$ be two integral preschemes, and $f:X\to Y$ a dominant morphism.
For every torsion-free quasi-coherent $\mathscr{O}_X$-module $\mathscr{F}$,
$f_*(\mathscr{F})$ is a torsion-free $\mathscr{O}_Y$-module.
\end{proposition}

\oldpage[I]{164}
\begin{proof}
Since $f_*$ is left exact
(\hyperref[0.4.2.1]{0, 4.2.1}), it suffices, by virtue of
(\hyperref[I.7.4.2]{7.4.2}), to prove the proposition when
$\mathscr{F}=(\mathscr{R}(X))^{(I)}$. Now, every non-empty open subset $U$ of
$Y$ contains the generic point of $Y$, hence $f^{-1}(U)$ contains the generic
point of $X$ (\hyperref[0.2.1.5]{0, 2.1.5}), so that we then have
\[
  \Gamma(U,f_*(\mathscr{F}))
  =\Gamma(f^{-1}(U),\mathscr{F})=(R(X))^{(I)}~;
\]
in other words, $f_*(\mathscr{F})$ is the simple sheaf with stalk
$(R(X))^{(I)}$, regarded as an $\mathscr{R}(Y)$-module, and it is evidently
torsion-free.
\end{proof}

\begin{proposition}[7.4.6]
\label{I.7.4.6}
Let $X$ be an integral prescheme and $x$ its generic point. For every
quasi-coherent $\mathscr{O}_X$-module of finite type $\mathscr{F}$, the
following conditions are equivalent: a) $\mathscr{F}$ is a torsion sheaf;
b) $\mathscr{F}_x=0$; c)
$\operatorname{Supp}(\mathscr{F})\neq X$.
\end{proposition}

By virtue of (\hyperref[I.7.3.5]{7.3.5}) and
(\hyperref[I.7.4.1]{7.4.1}), the relations $\mathscr{F}_x=0$ and
$\mathscr{F}\otimes_{\mathscr{O}_X}\mathscr{R}(X)=0$ are equivalent, hence
a) and b) are equivalent; on the other hand,
$\operatorname{Supp}(\mathscr{F})$ is closed in $X$
(\hyperref[0.5.2.2]{0, 5.2.2}), and since every non-empty open subset of
$X$ contains $x$, b) and c) are equivalent.

\begin{env}[7.4.7]
\label{I.7.4.7}
We extend (by abuse of language) the definitions of
(\hyperref[I.7.4.1]{7.4.1}) to the case where $X$ is a
\emph{reduced} prescheme having only a \emph{finite} number of irreducible
components; it then follows from (\hyperref[I.7.3.4]{7.3.4}) that the
equivalence of a) and c) in (\hyperref[I.7.4.6]{7.4.6}) is still valid for
such a prescheme, provided c) is replaced by the condition that
$\operatorname{Supp}(\mathscr{F})$ contains no irreducible component of $X$.
\end{env}
